\honest{The Second Law of Thermodynamics, and Engines of Cognition}{Eliezer Yudkowsky}{2008}

No machine of wheels and gears can create energy from nothing, I begin, because each part conserves energy, so the whole must too. If you claim that your machine does, you must name the wheel that breaks physics. I make this point three times: with particles, with wheels and gears, and with a drive that pushes without pushing anything back. It will come back in my last sentence.

A machine that turns warm water into electricity and ice is forbidden by the Second Law, which is ``essentially Bayesian in nature.'' I insist: ``Yes, really.'' \nb{This is E. T. Jaynes's view, published in 1957, and it is disputed among philosophers of physics.} Later I credit Jaynes for a proverb.

Next I explain phase space. The state of a whole system is one point in a very large space, and the laws of physics move the point. The volume of any region of that space stays the same over time. This is Liouville's theorem. I state it three times, twice before I name it, and each time I say it ``can be proven.'' I do not prove it.

A toy example follows. System X starts in one known state, and system Y in one of four. A minute later, Y is in one known state and X is in one of four. Uncertainty has moved from Y to X. If we do not bother to track exactly which four states X could be in, we say it could be in any of seven, and the total uncertainty appears to rise.

``The Second Law of Thermodynamics is a corollary of Liouville's Theorem,'' I say. In the next paragraph I say that the increase of entropy is an ``appearance,'' produced by drawing a simple boundary around a complicated region. I do not say which of the two I mean.

Then the thesis: the Second Law is a statement about your beliefs about a system. Suppose you learned the exact position and speed of every molecule in a glass of water. Is the water colder? ``Yes! Yes it is!'' \nb{Knowing where the molecules are does not slow any of them down, and a thermometer would read the same. What changes is how much work you could get out of the water.} I have redefined ``colder'' to mean ``more work available to someone who knows.'' I do not say so.

Maxwell's demon, a creature that lets fast molecules through a door and keeps slow ones back, cannot make a free refrigerator, I say, because it ``generates entropy in the process of inspecting the gas molecules.'' \nb{That was the standard account until 1982, when Charles Bennett showed that looking can in principle be free, and that the unavoidable cost comes from resetting the demon's memory.} In a parenthesis I call that distinction ``just a matter of words and perspective; the math is unambiguous.'' I show no math. My own table for the demon ends with its memory in one of four states, which is Bennett's account. Knowledge, I say, ``has to be represented in a brain, and that makes it as physical as anything else.'' So ``one subsystem cannot increase in mutual information with another subsystem, without (a) interacting with it and (b) doing thermodynamic work.''

After about 2,400 words of physics comes my conclusion, in bold: ``To form accurate beliefs about something, you really do have to observe it.'' \nb{That was true before thermodynamics.}

Thinking machines, I add, must give off waste heat like car engines, so ``cold rationality'' is true in a sense Hollywood never dreamed of.

Then the use of all this: ``unless you can tell me which specific step in your argument violates the laws of physics by giving you true knowledge of the unseen, don't expect me to believe that a big, elaborate clever argument can do it either.'' The physics has become a rule for debates.
